% allgem. Dokumentenformat
\documentclass[a4paper,12pt,headsepline]{scrartcl}

% deutsche Silbentrennung
\usepackage[ngerman]{babel}

% Deutsche Sonderzeichen benutzen 
\usepackage{ngerman}

% Eurozeichen einbinden
\usepackage[right]{eurosym}

% Umlaute unter UTF8 nutzen
\usepackage[utf8]{inputenc}

% Zeichenencoding
\usepackage[T1]{fontenc}

% Euro-Währung
\usepackage{eurosym}

%Variablen welche innerhalb der gesamten Arbeit zur Verfügung stehen sollen
\newcommand{\titleDocument}{Große Studienarbeit}
\newcommand{\subjectDocument}{im Studiengang Informationstechnik}

% weitere Pakete
% Grafiken aus PNG Dateien einbinden
\usepackage{graphicx}

\usepackage{lmodern}
\usepackage{fix-cm}

% floatende Bilder ermöglichen
%\usepackage{floatflt}

% mehrseitige Tabellen ermöglichen
\usepackage{longtable}

% Unterstützung für Schriftarten
%\newcommand{\changefont}[3]{ 
%\fontfamily{#1} \fontseries{#2} \fontshape{#3} \selectfont}

% Packet für Seitenrandabständex und Einstellung für Seitenränder
\usepackage{geometry}
\geometry{left=3.5cm, right=2cm, top=2.5cm, bottom=2cm}

% Paket für Boxen im Text
\usepackage{fancybox}

% bricht lange URLs "schoen" um
\usepackage[hyphens,obeyspaces,spaces]{url}

% Paket für Textfarben
\usepackage{color}

% Mathematische Symbole importieren
\usepackage{amssymb}

% auf jeder Seite eine Überschrift (alt, zentriert)
%\pagestyle{headings}

% erzeugt Inhaltsverzeichnis mit Querverweisen zu den Kapiteln (PDF Version)
\usepackage[bookmarksnumbered,pdftitle={\titleDocument},hyperfootnotes=false]{hyperref} 
%\hypersetup{colorlinks, citecolor=red, linkcolor=blue, urlcolor=black}
%\hypersetup{colorlinks, citecolor=black, linkcolor= black, urlcolor=black}

% neue Kopfzeilen mit fancypaket
\usepackage{fancyhdr} %Paket laden
\pagestyle{fancy} %eigener Seitenstil
\fancyhf{} %alle Kopf- und Fußzeilenfelder bereinigen
\fancyhead[L]{\nouppercase{\leftmark}} %Kopfzeile links
\fancyhead[C]{} %zentrierte Kopfzeile
\fancyhead[R]{\thepage} %Kopfzeile rechts
\renewcommand{\headrulewidth}{0.4pt} %obere Trennlinie
%\fancyfoot[C]{\thepage} %Seitennummer
%\renewcommand{\footrulewidth}{0.4pt} %untere Trennlinie

% für Tabellen
\usepackage{array}

% Runde Klammern für Zitate
%\usepackage[numbers,round]{natbib}

% Festlegung Art der Zitierung - Havardmethode: Abkuerzung Autor + Jahr
\bibliographystyle{alphadin}

% Schaltet den zusätzlichen Zwischenraum ab, den LaTeX normalerweise nach einem Satzzeichen einfügt.
\frenchspacing

% Paket für Zeilenabstand
\usepackage{setspace}

% für Bildbezeichner
\usepackage{capt-of}

% für Stichwortverzeichnis
\usepackage{makeidx}

% für Listings
\usepackage{listings}
\lstset{numbers=left, numberstyle=\tiny, numbersep=5pt, keywordstyle=\color{black}\bfseries, stringstyle=\ttfamily,showstringspaces=false,basicstyle=\footnotesize,captionpos=b}
\lstset{language=java}

% Indexerstellung
\makeindex

% Abkürzungsverzeichnis
\usepackage[german]{nomencl}
\let\abbrev\nomenclature

% Abkürzungsverzeichnis LiveTex Version
\renewcommand{\nomname}{Abkürzungsverzeichnis}
\setlength{\nomlabelwidth}{.25\hsize}
\renewcommand{\nomlabel}[1]{#1 \dotfill}
\setlength{\nomitemsep}{-\parsep}
\makenomenclature
%\makeglossary

% Abkürzungsverzeichnis TeTEX Version
% \usepackage[german]{nomencl}
% \makenomenclature
% %\makeglossary
% \renewcommand{\nomname}{Abkürzungsverzeichnis}
% \setlength{\nomlabelwidth}{.25\hsize}
% \renewcommand{\nomlabel}[1]{#1 \dotfill}
% \setlength{\nomitemsep}{-\parsep}

% Disable single lines at the start of a paragraph (Schusterjungen)
\clubpenalty = 10000
% Disable single lines at the end of a paragraph (Hurenkinder)
\widowpenalty = 10000
\displaywidowpenalty = 10000

\begin{document}
% hier werden die Trennvorschläge inkludiert
%hier müssen alle Wörter rein, welche Latex von sich auch nicht korrekt trennt bzw. bei denen man die genaue Trennung vorgeben möchte
\hyphenation{
Film-pro-du-zen-ten
Lux-em-burg
Soft-ware-bau-steins
zeit-in-ten-siv
}

%Schriftart Helvetica
%\changefont{phv}{m}{n}

% Leere Seite am Anfang
\newpage
\thispagestyle{empty} % erzeugt Seite ohne Kopf- / Fusszeile
\section*{ }

% Titelseite %
% das Papierformat zuerst
%\documentclass[a4paper, 11pt]{article}

% deutsche Silbentrennung
%\usepackage[ngerman]{babel}

% wegen deutschen Umlauten
%\usepackage[ansinew]{inputenc}

% hier beginnt das Dokument
%\begin{document}


\thispagestyle{empty}

%\begin{figure}[t]
% \includegraphics[width=0.6\textwidth]{abb/fh_koeln_logo}
%\end{figure}

\begin{figure}[t]
 \centering
 \includegraphics[width=0.6\textwidth]{abb/logo1}
\end{figure}


\begin{verbatim}


\end{verbatim}

\begin{center}
\Large{Duale Hochschule Baden-Württemberg}\\
\Large{- Karlsruhe -}\\
\end{center}


\begin{center}
\Large{Fakultät für Informatik}
\end{center}
\begin{verbatim}




\end{verbatim}
\begin{center}
\doublespacing
\textbf{\LARGE{\titleDocument}}\\
\singlespacing
\begin{verbatim}

\end{verbatim}
\textbf{{~\subjectDocument}}
\end{center}
\begin{verbatim}

\end{verbatim}
\begin{center}

\end{center}
\begin{verbatim}

\end{verbatim}
\begin{center}
\textbf{zur Erlangung des akademischen Grades \\ Bachelor of Engineering}
\end{center}
\begin{verbatim}






\end{verbatim}
\begin{flushleft}
\begin{tabular}{llll}
\textbf{Thema:} & & Social Engineering & \\
& & \\
\textbf{Autor:} & & Mario Philipp Waxenegger <mariowaxenegger@gmail.de>& \\
& & MatNr. 3981981 & \\
& & \\
\textbf{Version vom:} & & \today &\\
& & \\
\textbf{Betreuer:} & & Ralf Brune &\\
\end{tabular}
\end{flushleft}

% römische Numerierung
%\pagenumbering{arabic}

% 1.5 facher Zeilenabstand
\onehalfspacing

% Sperrvermerk
\section*{Sperrvermerk}
\textcolor{red}{
Die vorliegende Arbeit beinhaltet interne und vertrauliche Informationen der Firma <Firmenname>.
Die Weitergabe des Inhalts der Arbeit im Gesamten oder in Teilen sowie das Anfertigen
von Kopien oder Abschriften - auch in digitaler Form - sind grundsätzlich untersagt.
Ausnahmen bedürfen der schriftlichen Genehmigung der Firma <Firmenname>.
}

% Einleitung / Abstract
\section*{Zusammenfassung}


%\begin{verbatim}

%

%\end{verbatim}

\section*{Abstract}


% einfacher Zeilenabstand
\singlespacing

% Inhaltsverzeichnis anzeigen
\newpage
\tableofcontents

% das Abbildungsverzeichnis
%\newpage
% Abbildungsverzeichnis soll im Inhaltsverzeichnis auftauchen
\addcontentsline{toc}{section}{Abbildungsverzeichnis}
% Abbildungsverzeichnis endgueltig anzeigen
\listoffigures

% das Tabellenverzeichnis
%\newpage
% Abbildungsverzeichnis soll im Inhaltsverzeichnis auftauchen
\addcontentsline{toc}{section}{Tabellenverzeichnis}
% \fancyhead[L]{Abbildungsverzeichnis / Abkürzungsverzeichnis} %Kopfzeile links
% Abbildungsverzeichnis endgueltig anzeigen
\listoftables

%% WORKAROUND für Listings
%\makeatletter% --> De-TeX-FAQ
%\renewcommand*{\lstlistoflistings}{%
%  \begingroup
%    \if@twocolumn
%      \@restonecoltrue\onecolumn
%    \else
%      \@restonecolfalse
%    \fi
%    \lol@heading
%    \setlength{\parskip}{\z@}%
%    \setlength{\parindent}{\z@}%
%    \setlength{\parfillskip}{\z@ \@plus 1fil}%
%    \@starttoc{lol}%
%    \if@restonecol\twocolumn\fi
%  \endgroup
%}
%\makeatother% --> \makeatletter
% das Listingverzeichnis
%\newpage
% Listingverzeichnis soll im Inhaltsverzeichnis auftauchen
\addcontentsline{toc}{section}{Listingverzeichnis}
\fancyhead[L]{Abbildungs- / Tabellen- / Listingverzeichnis} %Kopfzeile links
\renewcommand{\lstlistlistingname}{Listingverzeichnis}
\lstlistoflistings
%%%%

% das Abkürzungsverzeichnis
%\newpage
% Abkürzungsverzeichnis soll im Inhaltsverzeichnis auftauchen
\addcontentsline{toc}{section}{Abkürzungsverzeichnis}
% das Abkürzungsverzeichnis entgültige Ausgeben
\fancyhead[L]{Abkürzungsverzeichnis} %Kopfzeile links
\nomenclature{UGC}{User Generated Content}
\nomenclature{CSS}{Cascading Style Sheets}
\nomenclature{JS}{JavaScript}
\nomenclature{SQL}{Structured Query Language}
\nomenclature{GPL}{GNU General Public License}
\nomenclature{GNU}{GNU is not Unix}
\nomenclature{LGPL}{GNU Lesser General Public License}
\nomenclature{XMPP}{Extensible Messaging and Presence Protocol}
\nomenclature{IM}{Instant Message}
\nomenclature{CMS}{Content Management System}
\nomenclature{RSS}{Really Simple Syndication}
\nomenclature{JSON}{JavaScript Object Notation}
\nomenclature{HTML}{Hypertext Markup Language}
\nomenclature{TDD}{Test-driven development}
\nomenclature{GUI}{Graphical User Interface}
\nomenclature{KPI}{Key Performance Indicator}
\nomenclature{WWW}{World Wide Web}
\nomenclature{OCR}{Optical Character Recognition}
\nomenclature{ERM}{Entity Relationship Modell}

\printnomenclature

% Definiert Stegbreite bei zweispaltigem Layout
\setlength{\columnsep}{25pt}

%%%%%%% EINLEITUNG %%%%%%%%%%%%
%\twocolumn
\newpage
\fancyhead[L]{\nouppercase{\leftmark}} %Kopfzeile links

% 1,5 facher Zeilenabstand
\onehalfspacing

% einzelne Kapitel
\section{Einleitung}\label{sec:2_einleitung}
\fachwort{Social Engineering} könnte man im Allgemeinen mit \ugspr{gesellschaftliches Ingeniuerwesen} übersetzen.
Jedoch trifft es die Bezeichnung \speech{angewandte Sozialwissenschaften} wesentlich besser.
Die Techniken hinter \fachwort{Social Engineering} reichen schon sehr weit in die Vergangenheit zurück, wohingegen der Begriff selbst erst seit einigen Jahren durch verschiedene Hacker wie z.B. \name{Kevin Mtnick} geprägt wurde.
Worum es sich bei \fachwort{Social Engineering} handelt soll im Rahmen dieser Ausführung in erster Linie geklärt werden.

Vielmehr stellt sich allerdings die Frage, wie \fachwort{Social Engineering} von der IT oder anderen Abteilungen gesehen und eingestuft wird.
Da es sich um einen recht neuen Terminus handelt, ist dieser womöglich noch nicht vollends in den Unternehmen und Einrichtungen angekommen.
Zwar kennt ein jeder die Rundmails eines nigerianischen Prinzen, der für sein Lebensglück noch schnell ein paar tausend Euro benötigt, jedoch handelt es sich bei solchen Maschen nur um die Spitze des Eisbergs \fachwort{Social Engineering} - die tatsächlichen Möglichkeiten und Gefahren reichen viel weiter.

Um dieser Frage nachzugehen wird das Thema sowohl aus der Sicht des Angreifers als auch aus der Sicht der potenziellen Zielperson analysiert und ausgearbeitet.
Dabei entpuppt sich \fachwort{Social Engineering} als Vielkampf der angewandten Wissenschaften.
Für einen Social Engineer gilt es oft die Bereiche Soziologie, Psychologie und einem gewissen Grad an schauspielerischem Können mit technischem Know-How zu verknüpfen.

\subsection{Ziel der Arbeit}\label{sec:ziel_der_arbeit}
Im Rahmen dieser Arbeit sollen einige Aspekte zum Thema Social Engineering überprüft und näher untersucht werden.
Der Kern der Ausarbeitung liegt im Ergebnis einer Umfrage.
Anhand dieser soll zunächst festgestellt werden, wie stark die Gefahren solcher Angriffe wahrgenommen werden.
Dabei gilt es außerdem herauszufinden, ob das Bewusstsein der Probanden für Social Engineering kontextabhängig ist.
Außerdem soll untersucht werden, ob es diesbezüglich Auffälligkeiten gibt, die darauf hindeuten, dass bestimmte Personengruppen stärker und weniger anfällig sind.
Ziel dieser Betrachtung ist es herauszufinden, ob bestimmte Personengruppen schlecht gegen Social Engineering Angriffe gewappnet sind, tatsächlich der Meinung sind für solche Angriffe keinen Nährboden bieten zu können.
Des weiteren sollen Korrelationen zwischen Anfälligkeit und anderen Attributen wie Alter, Geschlecht oder ähnlichem hergestellt werden.
Das Ziel dieser Ausarbeitung ist es, ein grundlegendes Verständnis von Social Engineering zu vermitteln
und anhand der Untersuchungen die Gefährlichkeit einzustufen. Zudem sollen in Anbetracht der Ergebnisse
Schutzmöglichkeiten eruiert werden.

Diese Ausarbeitung behandelt jedoch nicht die konkrete Vorgehensweise von Social Engineers.
D.h. es werden weder Tools analysiert noch werden detaillierte Fallanalysen durchgeführt.
Es werden lediglich dem Verständnis förderliche Beispiele erläutert.
Zudem sollen die beiden Kategorien telefonische und persönliche (direkte) Attacke nur theoretisch behandelt werden.
Da die Arbeit ebenso auf zahlreiche psychologische Grundlagen zurückgreift, werden Themen wie \fachwort{Nonverbale} Kommunikation ebenso kurz aufgegriffen.
 	
% keine genaue erklärung was in seminaren unterrichtet wird! lediglich orientierung an psych grundlagen
% keine Beschreibung technischer Grundlagen für Spoofing etc.

\subsection{Vorgehensweise}\label{sec:vorgehensweise}
Zu Beginn der Arbeit steht eine umfangreiche Grundlagenrecherche.
Dafür werden die Themen Vertrauen, Manipulation, Kommunikation, Risikomanagement und natürlich Social Engineering selbst analysiert.
Aus den Ergebnissen dieser Recherche wird zum einen ein Fragebogen entworfen, welcher die Empfänglichkeit für Social Enginneering überprüfen soll, und zum anderen wird eine Phishing-Mail konstruiert.
Dabei sollen aus der Theorie gewonnene Erkenntnisse vertieft und in der Anwendung geprüft werden.

Für die Recherche des Themas \fachwort{Social Engineering} wird Literatur von auf diesem Gebiet wegweisenden Experten zu Rate gezogen. Darunter fallen die Autoren \name{Kevin Mitnick}, \name{Christopher Hadnagy} und \name{Ian Mann}.
Das Thema \ugspr{Vertrauen}, bei dem es sich um ein psychologisches wie soziologisches Phänomen handelt, wird zu Großen Teilen anhand der Autoren \name{Niklas Luhmann} und \name{Bruce Schneier} aufgearbeitet.
Ausgangspunkt für die Recherche zum Thema Kommunikation und Kommunikationsmodelle bietet die Arbeit \speech{Grundlagen der Kommunikation} von \name{Markus Plate}, welche unter anderem auf die grundlegenden Ergebnisse von \name{Paul Watzlawick} aufbauen.
Das Werk \speech{Die Psychologie des Überzeugens} von \name{Robert Cialdini} sowie die Arbeiten von \name{Hadnagy} bieten die Basis des Kapitels zum Thema \fachwort{Manipulation}.
\fachwort{Risikomanagemt} wird anhand der Autoren \name{Dan Borge} und \name{Ian Mann} näher erläutert.

\subsection{Leitfaden}\label{sec:aufbau_der_arbeit}
\Kapitel{sec:social_engineering} stellt den Begriff Social Engineering und die Idee dahinter vor.
Dabei werden auch Beispiele für mögliche Angriffe gegeben.

In Kapitel \Kapitel{sec:risikoanalyse} wird der Begriff Risiko näher erläutert.
Da es sich bei Social Engineering ebenfalls um ein Risiko handelt werden die Analyse und die Handhabung
von Risiken sowie damit verbundene Schwierigkeiten genannt.

Um Social Engineering zu verstehen, ist es wichtig die zu Grunde liegenden Prinzipien zu erläutern.
In \Kapitel{sec:psychologische_grundlagen} werden einige Verhaltensmuster vorgestellt und aufgezeigt
wie diese gezielt ausgenutzt werden können.
Zentrale Themen sind hierbei vor allem Vertrauen, das Zusammenspiel von Bewusstsein und Unterbewusstsein.

Kapitel \Kapitel{sec:umfrageanalyse} befasst sich mit der Auswertung einer Umfrage, die im Rahmen der
Studienarbeit durchgeführt worden ist.

Im darauf folgenden \Kapitel{sec:praxis} wird die praktische Anwendung einer Social Engineering Attacke 


\section{Social Engineering}\label{sec:socialengineering}
Dieses Kapitel beschreibt den Begriff Social Engineering und grenzt die dazu gehörenden Techniken von üblichen Vorgehensweisen des Hackens ab. Dabei wird in Kapitel 2.1 zunächst ein Definition gegeben. Im darauf folgenden Kapitel wird die Einfachheit solche Angriffe gegenüber technischer Angriffe hervorgehoben. Kapitel 2.3 zählt die von Social Engineering herrührenden Risiken für Unternehmen mit kritischen Daten auf und bewertet diese. Zum Schluss (Kapitel 2.4) werden Beispiele für verschiedene Methoden gegeben.

\subsection{Definition}\label{sec:definition}
Um den Begriff Social Engineering korrekt einordnen zu können müssen zunächst herkömmliche Aspekte der Informationssicherheit betrachtet werden. Bei diesen handelt es sich zum einen um physikalische Zugriffskontrolle (z.B. Identitätsprüfungen an Türen)und zum anderen um IT-Sicherheit (meistens wird allerdings bei IT-Sicherheit lediglich von „Sicherheit“ gesprochen). Diese beiden Sicherheitsaspekte haben zweifellos ihren Platz und ihre Berechtigung. In vielen Fällen werden aber nur diese Art Angriffe berücksichtigt, denen solche Systeme entgegenarbeiten sollen. Offensichtlich ist, dass sich IT-Sicherheit und physikalische Sicherheit ausschließlich auf ein Unternehmen beschränken, welches lediglich aus IT-Systemen und Gebäuden mit Türen und Fenstern bestehen. Eine der essenziellsten Kernkomponenten des Unternehmens wird dabei gänzlich übersehen. Der Mitarbeiter stellt das größte Kapital des Unternehmens dar. Durch Social Engineering Techniken wird er jedoch zugleich zur größten Sicherheitslücke eines Unternehmens. Dies liegt größtenteils an den mangelnden Gegenmaßnahmen die es zu Social Engineering Attacken gibt.

Nun stellt sich weiter die Frage, was einen solchen Angriff ausmacht. Social Engineering Attacken zielen darauf ab, bestimmte Personen dahingehend zu manipulieren bestimmte Informationen herauszugeben oder Handlungen auszuführen, für die der EIngreifer selbst keine Berechtigung besitzt. Um solche Handlung auszulösen werden verschiedenste Techiken zur Täuschung herangezogen, die alle auf psychologischen Erkenntnissen beruhen. Auf diese wird in Kapitel 3 Psychologische Grundlagen tiefer eingegangen.

\subsection{Alternative zu technischen Methoden}\label{alternativezutechnischenmethoden}
In den letzten Jahren erfreuten sich Social Engineering Angriffe zunehmender Beliebtheit. Diese Entwicklung ist nicht ohne Grund zu beobachten. Während seit Beginn des Informationszeitalters auch die Angriffe auf Datenbestände immer häufiger und vor allem gefährlicher geworden sind, wurden entsprechend starke Gegenmaßnahmen entwickelt. Diese beschränken sich bis heute auf die Apsekte der physikalischen Sicherheit und der IT-Sicherheit. Auch heute noch spielt vor allem die IT-Sicherheit in vielen Unternehmen sicherlich gerechtfertigt eine übergeordnete Rolle. In jedem Unternehmen finden sehr wirkungsvolle aber auch ebenso kostenintensive Abwehrmechanismen ihre Anwendung. Diese stellen zwar kein unüberwindbares Hindernis dar, halten aber dennoch vielen Angriffen stand oder schrecken bereits vor einem Versuch ab. Während große Sicherheitsfirmen hierfür teure Software- und Hardwarelösungen anbieten, gibt es derzeit kein besonders großes Angebot an Abwehrmaßnahmen, die der Prävention von manipulativen Angriffen auf die Mitarbeiter dienen.

Es liegt also auf der Hand, dass immer mehr Angriffe nicht mehr herkömmlich auf die IT-Systeme direkt gerichtet werden, sondern entsprechende Mitarbeiter als Ziel haben. Der zu konventionellen Methoden gesparte Aufwand ist größer als er erwartet wird. Im nächsten Kapitel werden einige gängige Methoden und deren Effizienz vorgestellt.

\subsection{Gängige Angrifffe}\label{gangigeangriffe}
In diesem Kapitel werden exemplarisch einige Angriffstechniken vorgestellt und es wird grob geklärt, warum diese Techniken besonders große Erfolgschancen bieten. Dabei werden einige psychologische Grundlagen vorweggenommen, welche in Kapitel 3 detailliert beschrieben werden. Zum Verständnis der folgenden Beispiele sind diese Grundlagen nicht notwendig, jedoch werden nach der Lektüre des Kapitels Psychologische Grundlagen die Vorgänge hinter diesen Beispielen noch sehr viel deutlicher erkennbar sein.

\subsubsection{Heimarbeit und Helpdesks}\label{heimarbeitundhelpdesks}
Ein beliebtes Ziel für Social Engineering Angriffe stellen Mitarbeiter dar, die von zu Hause aus arbeiten. Dabei können die Mitarbeiter direkt als Ziel genommen werden. Mitarbeiter, die zu großen Teilen von zu Hause aus arbeiten, wissen oftmals nicht über alle Kollegen Bescheid und können auf Anrufe eines Angreifers, der sich als vermeintlicher Arbeitskollege ausgibt, vertrauliche Informationen herausgeben. Dabei wird Distanz zum Unternehmen als Schwachstelle ausgenutzt.

Dieses Angriffsziel ist allerdings zweischneidiger Natur. Auch in die andere Richtung können Angriffe vollzogen werden. Dabei wird der IT-Helpdesk als Zielscheibe gewählt. Mitarbeiter eines solchen Helpdesks sind oft geschult darauf freundlich und zuvorkommend zu handeln. Da gegenüber Heimarbeitern eine besonders große Hilfsbereitschaft an den Tag gelegt wird, ist es eine beliebte Taktik sich als Heimarbeiter auszugeben und so gewünschte Informationen zu erhalten.

\subsubsection{Neue Angestellte}\label{neueangestellte}
Dieses Szenario ist zwar etwas seltener anzutreffen aber durchaus nicht zu unterschätzen, da auf diese Weise bereits Zugriffsrechte auf bestimmte Bereiche des Unternehmens oder der Datenbasis gewährt wird. Auch in solchen Situationen wird die Hilfsbereitschaft von anderen Mitarbeitern ausgenutzt, denn gerade neue Mitarbeiter benötigen am Anfang viel Hilfe. Außerdem bringt man neuen Mitarbeitern mehr Nachsehen entgegen.

In einer solchen Position ist es für einen Angreifer ein leichtes an eine Vielzahl wichtiger Informationen zu kommen oder zusätzliche Schwachstellen ausfindig zu machen.

Wenn neue Mitarbeiter nach kurzer Zeit ohne plausiblen Grund Kündigen, ist dies oft ein Indiz dafür, dass es sich um einen Angriff gehandelt haben kann. Um solche Situationen zu verhindern, ist es angebracht für neue Mitarbeiter ausreichende Background-Checks vorzunehmen.

\subsubsection{Angriff auf Mitwisser}\label{angriffaufmitwisser}
Manchmal stellt es den Angreifer vor große Probleme das Ziel direkt zu attakieren, da evtl. keine brauchbaren Sicherheitslücken ausfindig gemacht werden können oder das Risiko eines direkten Angriffes zu hoch wäre. Hier bietet es sich für den Angreifer an sich ein kooperierendes Unternehmen als Ziel zu wählen, welches größere Sicherheitslücken aufweist. Die Informationen die Dritte über das Zielunternehmen haben, können entscheidend für den tatsächlichen Angriff bieten. Auf ein solches drittes Unternehmen können dann Methoden wie bereits beschrieben angewendet werden.

\subsubsection{Kombination sozialer und technischer Angriffe}\label{kombination}
 Oft ist mit rein manipulativen Mitteln das gewünschte Ziel nicht erreichbar. Dabei nutzen Angreifer häufig eine Mischform von Social Engineering Techniken und herkömmlichen Angriffstechniken. Dafür kann ein Angreifer beispielsweise einen mit Malware infizierten USB-Stick in einer Büroetage liegen lassen. Die Neugier einiger Mitarbeiter wird dabei größer sein als die Vernunft und es ist keine seltenheit, dass ein solcher USB-Stick den Weg in die Buchse eines Mitarbeiter-PCs findet.

EIne sehr bekannte technsiche Methode, die sich Social Engineering Techniken bedient, ist das Senden von Phishing-Mails. Dabei werden an eine große Menge von Zielen E-Mails versendet, die sich verschiedene menschlichen Schwächen zu Nutze machen wie z.B. Leichtgläubigkeit, Kurzsichtigkeit oder Gier. Diese Technik findet allerdings weniger Anwendung bei Angriffen auf Unternehmen. Der Ausbeute eines solchen Angriffs ergibt sich nicht aus einem einzigen Angriff sondern vielmehr aus der Menge. Tatsächlich ist es eine Minderheit im Promillbereich, die auf solche Phishing-Mails reagiert. Sendet der Angreifer 1000000 Mails aus, bei denen pro Aktion \EUR{1000} erbeutet werden können, genügt eine Erfolgsrate von 0,01\% um 100 erfolgreiche Angriffe verbuchen zu können, was einem Gewinn von \EUR{100000} entspricht.

\section{Risikoanalyse und -management}\label{sec:risikoanalyse}
In den vorigen Kapiteln ist das Prinzip von Social Engineering erläutert worden und somit klar
geworden wie leicht solche Angriffe von statten gehen können und welche Gefahren im Vergleich zu
herkömmlichen Angriffsmethoden bestehen.
Dennoch wird das von Social Engineering herrührende Risiko oft unterschätzt.
Um diesen Gefahren zielgerichtet entgegenzuwirken ist es wichtig ein grundlegendes, allgemeines
Verständnis von Risikofaktoren, deren Bewertungen und der Handhabung mit ihnen zu bekommen.
Dieses Kapitel vermittelt zunächst wobei es sich bei dem Begriff Risiko im Allgemeinen handelt.
Es wird außerdem geschildert wie man mit bereits erkannten Risiken umgehen kann (Kapitel 3.1.1) und
aus welchen Gründen Risiken meist kein angemessenes Risikomanagement seine Anwendung findet.
Abschließend werden diese allgemeinen Erkenntnisse auf den Bereich der speziellen Risikofaktoren durch
Social Engineering Angriffe übertragen (Kapitel 3.2).

\subsection{Allgemeines}\label{sec:allgemeines}
Risiken im Allgemeinen und ihre Handhabung treten in nahezu allen Situationen des täglichen Lebens
auf.
Dabei kann es sich um solche Situationen trivialer Natur handeln wie der Wahl eines Getränks bis hin
zur Entscheidung für oder gegen die Investition einer Aktie.
Auch wenn beide Situationen sehr unterschiedlich anmuten, haben sie doch sehr viel gemeinsam.
Sowohl die Wahl eines Getränks als auch die Investition in eine Aktie bilden ein gewisses Risiko.
Offensichtlich bieten die beiden Aktionen unterschiedliche große Risiken, weswegen es notwendig ist
solche Risiken korrekt einzustufen und zu bewerten.
Ein wichtiges Merkmal für ein Risiko ist, dass ein jedes Risiko mit einer Entscheidung verknüpft ist,
wobei es sich dabei auch um die Entscheidung nichts zu tun handeln kann.

In Bezug auf Informationssicherheit kann eine Gefahr bedeuten, dass geschäftskritische Daten für
nicht-autorisierte Nutzer zugänglich gemacht werden.
Trifft man in einem solchen Fall die Entscheidung nichts zu tun, ist das damit verbundene Risiko
entsprechend hoch.

Die Auseinandersetzung mit Entscheidungen, den damit Verbunden Risiken und den Möglichkeiten des
zweckmäßigen Umgangs ist vor allem wertvoll für die positive Beeinflussung zukünftiger Situationen.
Die Analyse von Risiken sollte nicht dazu verwendet werden Ereignisse zu erklären welche der
Vergangenheit angehören.

In Bezug auf den Umgang mit Risiken gibt es verschieden charakterisierte Herangehensweisen, welche
freilich nicht in der Reinform auftreten, aber generell oft wiederzuerkennen sind.

Die fatalistische Herangehensweise zeichnet sich besonders durch die damit verbundene Passivität aus.
Dabei wird davon ausgegangen, dass zukünftige Ereignisse weder abwendbar noch voraussehbar sind.
Auffällig ist, dass das Handeln ausschließlich aus Reaktionen besteht.

Ein anderes Extrem ist der fanatische Ansatz. Hierbei wird von einer konkreten Zukunft ausgegangen,
welche man in diesem Zusammenhang als Vision bezeichnen kann.
Abweichungen von dieser Vision in Form anderer Möglichkeiten werden gänzlich ignoriert und lösen auch
keine notwendigen Reaktionen aus.

Dass diese beiden Handlungsweisen nicht zielführend sind, ist offensichtlich.
Allerdings ist auch eine rein wissenschaftliche Beurteilung von Risikosituationen nicht von nutzen.
Ein rein wissenschaftlicher Ansatz ist zwar unvoreingenommen, zeichnet sich aber durch die
Notwendigkeit einer stichhaltigen Beweislage aus.
Da diese für in der Zukunft liegende Ereignisse nicht vorhanden ist, würde der rein wissenschaftliche
Ansatz nie zu einer endgültigen Entscheidung führen.

Der pragmatische Ansatz geht von der Annahme aus, dass die Zukunft zwar ungewiss ist, aber nicht
gänzlich unvorhersehbar ist.
Dabei sollen die Chancen positiv verändert werden, während eine mangelnde Beweislage akzeptiert wird.

Dan Borge S. 1-8, 21-39

\subsection{Risikomanagement}\label{sub:risikomanagement}
Der Umgang mit Risiken setzt sich aus mehreren Schritten zusammen.
Der erste aber auch schwierigste Schritt ist die Identifizierung von Risiken.
Die Schwierigkeit besteht im Grunde darin, dass nicht klar ist, wonach überhaupt gesucht wird.
Das liegt an der bislang noch spärlichen Kategorisierung von Risiken.
In der Informationssicherheit sind jedoch viele Risiken bereits bekannt, obgleich sich die Methoden
schnell weiterentwickeln.
Ein Risiko stellt auch die Bedrohung durch Social Engineering Angriffe dar.

Gängige Werke zum Thema Risikomanagement empfehlen beim Umgang mit Risiken verschiedene Taktiken.
Dabei wird beispielsweise die Vermeidung, das Verkaufen und Verteilen der Risiken vorgeschlagen.
Im Falle von gespeicherten Informationen stellen diese Lösungsansätze allerdings keine brauchbaren
Alternativen dar.
Zwei Möglichkeiten für den Umgang mit Risiken solcher Art sind zum einen die Versicherung für den
Fall, dass ein Schadensfall eintritt, zum anderen die gezielte Absicherung.
Da mit dem Verlust solcher Daten an Dritte nicht nur finanzieller Schaden sondern auch ein immenser
Imageschaden einhergeht, ist auch die Herangehensweise mit einer Versicherung nicht zweckgemäß.
In aller Regel bildet nur ein gezielter Schutz vor Angriffen eine angemessene Art des
Risikomanagements.

// to do: Vorgehensweise für Risk-Management in IT-Sec recherchieren

Dan Borge S. 45-58

\subsection{Ursachen für schlechtes Risikomanagement}\label{sec:ursachen}
Oft stehen dem korrekten Umgang mit Risiken allerdings einige Hürden im Weg. Zum einen sind es sicherlich Schätzungen und Vereinfachungen, die eine Fehlerquelle bilden. Viel gefährlicher sind jedoch solche Fehlerquellen menschlicher Natur. Diese Hindernisse zur rationalen Entscheidungsfindung gründen auf grundlegenden Eigenschaften der menschlichen Psyche (welche wiederum die Grundlage für die Techniken des Social Engineering bildet).

Menschen neigen grundsätzlich dazu die eigenen Möglichkeiten zu überschätzen.
Möglichkeiten, denen eine sehr geringe Wahrscheinlichkeit zugeordnet ist, werden oft ausgeblendet und
als unmöglich bezeichnet.
Die Tendenz bei Menschen geht allerdings dahin, dass dies stärker bei Risiken geschieht als bei der
Wahrnehmung von Möglichkeiten.
Die Ursache hierfür ist das Phänomen des menschlichen Optimismus.
Negative Resultate werden vielmals unterschätzt und unverhältnismäßig hohe Risiken werden für Chancen
mit geringer Aussicht auf Erfolg aufgenommen.

Des Weiteren ist oft eine fälschliche Betrachtung vergangener Ereignisse zu beobachten.
Im Detail bedeutet das, dass eingetretenen Ereignissen der Vergangenheit im Nachhinein nachgesagt
wird, sie seien absehbar gewesen, auch wenn sie vor ihrem Eintreten als unwahrscheinlich eingeordnet
worden sind.
Mit dieser rückblickenden Betrachtung werden aus optimistischen Herangehensweise resultierende
Fehlentscheidungen gerechtfertigt und beschönigt.

Vielen Menschen fällt es außerdem schwer an die Beliebigkeit und die Zufälligkeit von Ereignissen zu
glauben.
Tatsächlich neigen sie dazu nach Mustern zu suchen und Ereignisse in eine Ordnung zu bringen.
So werden beispielsweise aus einer Entwicklung über eine relativ kurze Zeit Rückschlüsse für die Zukunft geschlossen.
Ähnlich dem Suchen von Mustern geht die Kurzsichtigkeit bei der Beurteilung von Ereignissen einher.
Dies ist unter anderem bei Fußballtrainern zu beobachten.
Schon eine kurze Serie an Niederlagen, verleitet Fans und Vereinsvorstand die Ursache beim Trainer zu
suchen.
Viele Trainer werden schon nach kurzen Zeiträumen wieder entlassen, ohne dass eine solide
Informationsbasis zur rationalen Entscheidungsfindung existiert.

Einen Fußballtrainer nach einer Serie von Niederlagen zu entlassen, ist zwar eine harte
Entscheidungsfindung, aber aus risikoanalytischen Gesichtspunkten immer noch sinnvoller als die
Trägheit des menschlichen Verhaltens.
Die Ursache hierfür liegt darin, dass Menschen eine falsche Aktion negativer bewerten als das
Unterlassen einer Aktion, auch wenn die Passivität das schlechtere Ergebnis liefert.
Gerade in Bezug auf die Informationssicherheit ist es undenkbar keine Gegenmaßnahmen zu treffen oder
im Falle eines Informationsverlustes in Tatenlosigkeit zu verharren.

Eine Bekannte Ursache für eine Fehleinschätzung von Risiken sind ironischerweise
Sicherheitsmechanismen wie der Sicherheitsgurt im Auto oder ein verbesserter Algorithmus zur
Berechnung des Risikos für Finanzanlagen.
Eben diese risikominimierenden Algorithmen sorgen dafür, dass sich der Anwender in Sicherheit fühlt
und damit in Summe ein größeres Risiko eingeht, als er ohne diese Absicherung eingegangen wäre.

Der letzte signifikante Punkt ist die Selbstzufriedenheit des Menschen. Darunter ist zu verstehen,
dass dem Menschen bekannte Risiken weniger gefährlich erscheinen als solche, die über seine Gewohnheit
hinausgehen.
Gerade in Bezug auf Social Engineering ist dieser Punkt erwähnenswert, denn es ist dem Menschen in der
Tat eine Gewohnheit mit anderen Menschen zu Interagieren.
So fällt es Menschen schwer das Risiko zu erkennen, welches von einem Gespräch ausgehen kann.
Grundsätzliche Skepsis ist allerdings ebenso kontraproduktiv und eine es muss abgeschätzt werden
inwiefern Sicherheitsvorkehrungen bei internen und externen Gesprächen eines Unternehmens sinnvoll
sind.
Das Finden dieser Balance wird im Kapitel 10 Gegenmaßnahmen genauer erörtert.


// to do: weitere Recherche Tversky, Kahnemann

Dan Borge S. 61-72

\subsection{Social Engineering als Risikofaktor}\label{sec:risikofaktor}
Bis zu diesem Kapitel wurde das Phänomen Social Engineering und die grundlegende Idee dahinter
beschrieben sowie der Umgang mit Risikofaktoren und den Hindernissen der menschlichen Psyche.

Social Engineering bildet unbestreitbar einen beachtenswerten Risikofaktor.
Dabei gibt es eine Vielzahl an möglichen Szenarien, die beobachtet und analysiert werden können.
Leider ist es zu oft der Fall, dass erfolgreiche Angriffe die auf Social Engineering basieren in den
seltensteten Fällen als solche identifiziert werden.
Gerade weil Angriffe dieser Art oft in Kombination mit technischen Methoden angewandt werden, wird ein
rein technischer Angriff dokumentiert.
Aus diesem Grund und eines allgemein fahrlässigen Risikomanagements wegen wird in vielen Unternehmen
die Gefahr von Social Engineering nicht angemessen ernst genommen.

%....

\section{Diskussion und Ausblick}\label{diskussion-und-ausblick}

Im Verlauf der vorliegenden Arbeit wurde zunächst an das Thema Social Engineering herangeführt und einige grundlegende Angriffsmuster beschrieben.
Die mit solchen Angriffen erreichte Effizienz ist erschreckend und ist ein klares Zeichen für Unternehmen diese Bedrohung ernst zu nehmen.
Vor allem die diffizile forensische Aufklärung solcher Angriffe macht es schwer genaue Zahlen zu erfassen, denn oftmals stehen hinter scheinbar rein technischen Angriffen die Methoden eines Social Engineers am Anfang der Kausalkette.
Zwar ist in sehr großen Konzernen die Gefahr bereits erkannt und korrekt eingeschätzt, jedoch ist es vor allem in kleinen und mittelständischen Unternehmen ein noch nicht ernst genug genommenes Thema.
Hacking Angriffe zielen zwar meist auf militärische Einrichtungen, Banken oder ähnliche Institutionen ab, jedoch liegt das Risiko für andere Branchen bei weitem nicht bei null.
Wie das Kapitel über \nameref{sec:risikoanalyse} gezeigt hat, werden viele Risiken wie Social Engineering oft nicht ernst genug genommen und zu oft darauf gehofft, dass ein negatives Ergebnis ausbleibt.
Gewissenhafte Datenschutzbeauftragte in Firmen sollten demnach überlegen das Unternehmen und seine Mitarbeiter gegen Angriffe dieser Art zu schützen.
Problematisch können an dieser Stelle selbstverständlich die aufkommenden Kosten für eine solide Implementierung von Abwehrmaßnahmen sein.
Es kann jedoch aktiv ein Verständnis der Vertraulichkeit von Daten gelebt werden, was schon mit wenigen Mitteln stark verbessert werden kann.

Die Umfrage bestätigt in einigen Punkten die Dringlichkeit einer intensiveren Aufklärung.
Zwar gaben nur wenige Probanden an einen Mitarbeiter mit defekten/fehlenden ID-Chip einzulassen, jedoch wurde auf Bitten anderer Szenarien weitaus freizügiger reagiert.
Zwar gibt es leichte Schwankungen zwischen Berufsgruppen welche in der EDV angesiedelt sind und anderen Berufsgruppen, diese beziehen sich jedoch größtenteils auf Situationen in denen es konkret um Aktivitäten an IT-Geräten geht.
Wenn es sich um Zugangskontrollen oder ähnliches handelt gleichen sich die Ergebnisse an.
Sicherlich bietet das Grundprinzip solcher Umfragen noch weiteres Potenzial zur genaueren Untersuchung.
Einzelne Zielgruppen können genauer analysiert werden, um somit genauere Vorstellungen von den Schwachstellen aufzudecken.



\section{Fazit}\label{fazit}

\onecolumn
% einfacher Zeilenabstand
\singlespacing
% Literaturliste soll im Inhaltsverzeichnis auftauchen
\newpage
\addcontentsline{toc}{section}{Literaturverzeichnis}
% Literaturverzeichnis anzeigen
\renewcommand\refname{Literaturverzeichnis}
\bibliography{Hauptdatei}

%% Index soll Stichwortverzeichnis heissen
% \newpage
% % Stichwortverzeichnis soll im Inhaltsverzeichnis auftauchen
% \addcontentsline{toc}{section}{Stichwortverzeichnis}
% \renewcommand{\indexname}{Stichwortverzeichnis}
% % Stichwortverzeichnis endgueltig anzeigen
% \printindex

\onehalfspacing
% evtl. Anhang
\newpage
\addcontentsline{toc}{section}{Anhang}
\fancyhead[L]{Anhang} %Kopfzeile links
\subsection*{Anhang}\label{anhang}


% Eidesstattliche Erklärung
\addcontentsline{toc}{section}{Eidesstattliche Erklärung}
\section*{Eidesstattliche Erklärung}
\thispagestyle{empty}

\begin{verbatim}

\end{verbatim}

\begin{LARGE}Eidesstattliche Erklärung zur <-Arbeit>\end{LARGE}
\begin{verbatim}


\end{verbatim}
Ich versichere, die von mir vorgelegte Arbeit selbstständig verfasst zu haben. Alle Stellen, die wörtlich oder sinngemäß aus veröffentlichten oder nicht veröffentlichten Arbeiten anderer entnommen sind, habe ich als entnommen kenntlich gemacht. Sämtliche Quellen und Hilfsmittel, die ich für die Arbeit benutzt habe, sind angegeben. Die Arbeit hat mit gleichem Inhalt bzw. in wesentlichen Teilen noch keiner anderen Prüfungsbehörde vorgelegen.



\begin{displaymath}
% use packages: array
\begin{array}{ll}
Unterschrift:~~~~~~~~~~~~~~~~~~~~~~~~~~~~~~~~~~~~~~~~~~
& Ort, Datum:~~~~~~~~~~~~~~~~~~~~~~~~~~~~~~~~~~~~~~~~~~
\end{array}
\end{displaymath}


% leere Abschlussseite
\newpage
\thispagestyle{empty} % erzeugt Seite ohne Kopf- / Fusszeile
\section*{ }

\end{document}